\documentclass[11pt]{article}

\usepackage[utf8]{inputenc}
\usepackage[T1]{fontenc}
\usepackage[ngerman]{babel}
\usepackage{amsmath,amssymb}
\usepackage[a4paper,margin=2.5cm]{geometry}
\usepackage{enumitem}
\usepackage[hidelinks]{hyperref}

\emergencystretch=5em

\newcommand{\code}[1]{\texttt{#1}}
\newcommand{\paperfile}{\path{LaTeX/rational_cut_and_project_gap_periods.tex}}
\newcommand{\researchfile}{\path{support/research.tex}}

\title{Codex-Review der aktuellen uncommitted Änderungen\\
\large Hinweise zu \paperfile{} und \researchfile{}}
\author{}
\date{31. Mai 2026}

\begin{document}
\maketitle

\section*{Kurzurteil}

Die aktuellen uncommitted Änderungen an \paperfile{} und \researchfile{} sind
im Kern plausibel und stärken die Recherche.  Ich sehe keinen mathematischen
Blocker.  Vor einem Commit würde ich aber zwei kleine Präzisierungen am
Paper-Text vornehmen und die neue stärkere Aussage im Rechercheprotokoll nur
dann so stehen lassen, wenn die direkte Konsultation beider Baake--Grimm-Bände
tatsächlich auf Volltext- bzw.\ Indexbasis erfolgt ist.

\section*{Findings}

\begin{enumerate}[leftmargin=2em,itemsep=8pt]
  \item \textbf{Mehrdeutiger Bezug von ``this literature'' im Paper.}

        In \paperfile{} um Zeile~1913 ist die neue Aussage zu
        \(D=\alpha^2+\beta^2\) inhaltlich plausibel:
        \(a^2+b^2\) tritt in Baake--Grimm Vol.~2 im Kontext von
        similar sublattices und coincidence-site lattices als
        Gaussian-integer norm form auf, also als enumerative Größe und nicht
        als Periodenmodulus einer Gap-Folge.

        Die Formulierung
        \begin{quote}
        The constant \(D = \alpha^2+\beta^2\) itself does appear in this
        literature, but in an enumerative role \dots
        \end{quote}
        ist jedoch leicht mehrdeutig, weil direkt davor
        Gähler--Hunton--Kellendonk erwähnt werden.  Die Norm-Aussage gehört
        nicht zu diesem Paper, sondern zu Baake--Grimm Vol.~2 /
        Baake--Zeiner.

        Bessere Formulierung:
        \begin{quote}
        In Baake--Zeiner's chapter in Vol.~2 of
        \code{\textbackslash cite\{BaakeGrimm2017\}}, the constant
        \(D=\alpha^2+\beta^2\) appears in an enumerative role---as the
        Gaussian-integer norm form used to count similar sublattices and
        coincidence-site lattices of \(\mathbb Z^2\)---rather than as the
        period modulus of a gap sequence.
        \end{quote}

  \item \textbf{Spezifische Kapitelreferenz wäre besser als nur der Sammelband.}

        Die neue Paper-Passage zitiert derzeit nur
        \code{\textbackslash cite\{BaakeGrimm2017\}}.  Für die spezifische
        Aussage über \(a^2+b^2\), similar sublattices und CSLs wäre eine
        Kapitelangabe robuster.  Das einschlägige Kapitel ist Baake--Zeiner,
        Ch.~3 von \emph{Aperiodic Order. Vol.~2}, mit dem Thema geometric
        enumeration problems for lattices and embedded \(\mathbb Z\)-modules.

        Minimaler Fix ohne neuen Bibitem:
        \begin{quote}
        \dots in Baake--Zeiner's Ch.~3 of \code{\textbackslash cite\{BaakeGrimm2017\}} \dots
        \end{quote}

        Noch sauberer wäre ein eigener Bibitem für das Kapitel, falls der
        Abschnitt im Paper bibliographisch sehr präzise sein soll.  Für die
        jetzige Related-Work-Abgrenzung reicht die Kapitelangabe im Text aber
        wahrscheinlich aus.

  \item \textbf{Stärkere Baake--Grimm-Aussage im Rechercheprotokoll nur mit
        tatsächlicher Volltext-/Indexgrundlage.}

        In \researchfile{} um Zeile~350 steht nun:
        \begin{quote}
        Both \emph{Aperiodic Order} volumes were consulted directly (not merely
        their frontmatter) \dots
        \end{quote}
        und der alte Residualpunkt zum offenen Baake--Grimm-Vol.-1-Indexcheck
        wurde entfernt.

        Das ist konsistent, \emph{wenn} wirklich beide Bände direkt konsultiert
        wurden.  Die gestagete \code{.gitignore}-Änderung für
        \code{support/bibliography/v1.pdf} und
        \code{support/bibliography/v2.pdf} spricht dafür.  Da diese PDFs nicht
        getrackt sind, ist diese Evidenz aber nur lokal.  Inhaltlich ist das
        okay; beim Commit sollten \code{.gitignore} und die Rechercheänderung
        zusammen betrachtet werden.

        Falls nur Frontmatter/Snippets vorlagen, sollte die Formulierung
        abgeschwächt werden, etwa:
        \begin{quote}
        The available full-text/index material for both volumes was consulted
        directly \dots
        \end{quote}
\end{enumerate}

\section*{Positive Punkte}

\begin{itemize}[leftmargin=2em,itemsep=4pt]
  \item Das Entfernen der alten Restlücke
        ``Baake--Grimm Vol.~1, verbatim index check'' ist konsistent, falls die
        direkte Konsultation tatsächlich erfolgt ist.
  \item Der neue Appendix \code{Where the fingerprint terms occur in
        Baake--Grimm} in \researchfile{} macht die Recherche nachvollziehbarer.
  \item Die Abgrenzung ``counting quantity, not a period modulus'' passt gut
        zur mathematischen Rolle von \(a^2+b^2\) in der CSL-/similar-sublattice-
        Literatur.
\end{itemize}

\section*{Build- und Statuscheck}

\begin{itemize}[leftmargin=2em,itemsep=4pt]
  \item \code{latexmk} für \paperfile{} läuft durch.
        Keine undefinierten Referenzen oder Zitate; nur die bekannten zwei
        Underfull-Box-Warnungen.
  \item \code{latexmk} für \researchfile{} läuft durch.
        Keine undefinierten Referenzen oder Zitate; nur ältere Layoutwarnungen
        (Overfull/Underfull boxes).
  \item \code{git diff --check} für die beiden Dateien meldet keine
        Whitespace-Probleme.
\end{itemize}

\section*{Empfehlung}

Ich würde die Paper-Passage minimal umformulieren, so dass sie explizit
Baake--Zeiner / Vol.~2 nennt, statt ``this literature'' zu verwenden.  Danach
sind die Änderungen aus meiner Sicht commit-reif, vorausgesetzt die direkte
Baake--Grimm-Konsultation in \researchfile{} entspricht dem tatsächlich
vorliegenden Arbeitsstand.

\end{document}
